
Geben Sie das Prüfungsschema eines Schadenregulierer wieder

0. Sachverhaltsaufklärung
1. Besteht überhaupt Deckung?
2. Ist die Haftung dem Grunde nach gegeben?
3. Welche Anspruchsgrundlage ist einschlägig?
4. Besteht evtl. ein Mitverschulden?


Geben Sie das Prüfungsschema für Schmerzensgeld wieder 


I. Anspruchsentstehung
1. Verletzung eines der in § 253 (Immaterieller Schaden) Abs. 2 BGB genannten Rechtsgüter 
(Körpers, der Gesundheit, der Freiheit oder der sexuellen Selbstbestimmung)
2. Schadenersatzanspruch dem Grunde nach wegen Verletzung
3. Höhe: Billige Entschädigung in Geld (Typischerweise Schmerzensgeldtabellen)
II.
Rechtsvernichtende Einwendungen
III.
Durchsetzbarkeit

Dies gilt auch für psych. Erkrankungen zentral dabei ist
- Primärverletzung
- Kausalität zwischen Schadenereignis und psychischer Erkrankung (§ 286 ZPO)
- Nicht nur allgemeines Lebensrisiko
- Keine Bagatelle
- Keine Begehrensneurose

Als Basis dient insbesondere (§ 843 BGB) Geldrente oder Kapitalabfindung:
(1) Wird infolge einer Verletzung des Körpers oder der Gesundheit die Erwerbsfähigkeit des Verletzten aufgehoben oder gemindert 
  oder tritt eine Vermehrung seiner Bedürfnisse ein, so ist dem Verletzten durch Entrichtung einer Geldrente Schadensersatz zu leisten.

(2) Auf die Rente findet die Vorschrift des § 760 Anwendung. Ob, 
  in welcher Art und für welchen Betrag der Ersatzpflichtige Sicherheit zu leisten hat, bestimmt sich nach den Umständen.
(3) Statt der Rente kann der Verletzte eine Abfindung in Kapital verlangen, wenn ein wichtiger Grund vorliegt.
(4) Der Anspruch wird nicht dadurch ausgeschlossen, dass ein anderer dem Verletzten Unterhalt zu gewähren hat.


Was sind Vermehrte Bedürfnisse?

Sie ergeben sich aus (§ 843 I BGB) und kennzeichen sich durch 
- unfallbedingte
- ständig wiederkehrenden Mehraufwendungen,
Beispiele für solche vermehrten Bedürfnisse sind z.B. 
- Umbaukosten für ein behindertengerechtes Fahrzeug oder eine behindertengerechte Wohnung, 
- Fahrtkosten für Arztbesuche, 
- Körperversorgung und Körperpflege, 
- Kosten für besondere Kleidung, 
- Lese- und Schreibhilfen,

Aber auch Haushaltsführungsschäden
diese werden beispielhaft berechnet wie folgt

Zeitaufwand wöchentlich x Grad der Behinderung in \% x Stundenlohn netto ./. kongruente
Leistungen Dritter x Haftungsquote = Schadenersatz

Sie sind abzugrenzen von

- Heilbehandlungskosten
- Allgemeine Lebenshaltungskosten



Wie analysieren Sie Erwerbsschäden?

1. Regress des Arbeitgebers
2. Vorteilsausgleich, Schadenminderung/ Umschulung
3. Fiktive Bestimmung des Erwerbsschadens

Die folgenden Zahlungen 
- Lohnfortzahlung
- Krankengeld
- Reha
werden dann im Regress vom 
- Arbeitgeber
- Gesetzliche Krankenkasse
- Regress Deutsche Rentenversicherung
an die Versicherung gestellt.

Im Falle dass die Arbeit nicht mehr in der Zukunft ausgeübt werden kann gibt es eine
grundsätzliche Verpflichtung zur Umschulung, erforderlichenfalls sogar Wechsel des Wohnortes
die Beweislast liegt beim Geschädigten bei haftungsausfüllende Kausalität ansonsten Strengbeweis.

Wird bei allen relevant, die noch nicht im Erwerbsleben stehen, Kinder, Auszubildende, Schüler,
Studenten
- Verzögerter Einstieg ist zu ersetzen
Vergleich von Ist- und Sollverlauf (Beweislast für Sollverlauf hat Geschädigter)
- je weiter weg der Geschädigte vom Erwerbsleben ist, desto schwerer/fiktiver die Prognose

Für welche Heilbehandlungskosten ist zu decken?

Tatsächlich entstandene, angemessene Kosten aller unfallbedingten (Kausalität) und erforderlichen
Heilbehandlungsmaßnahmen.

Als Spezialfall der privaten Krankenver.
trägt ein Verletzter die Heilbehandlungskosten selbst, indem er sich einer Privatbehandlung
unterzieht, so bleibt er hinsichtlich des Schadenersatzes aktiv legitimiert.
▪ ein Übergang auf die Krankenkasse findet nicht statt


Was sind Vor/Nachteile einer Kapitalabfindung

Vorteile VR:
- Effizienzgründe
- Wirtschaftliche Ersparnis
- "Veränderungsrisiko“ eher auf Anspruchstellerseite
Nachteil VR und VN
- Vorversterblichkeitsrisiko
- Doppelzahlungen
- Minderjährigkeit und Gefahr des „Verprassens“ durch Eltern?

Direktgeschädigter
- Anspruch auf Abfindung in Kapital nur, wenn „wichtiger Grund“ i.S.d. § 843 Abs. 3 BGB
- Ansonsten kein Anspruch auf Kapitalisierung, sondern Verhandlungssache
- Prozessrisiko abwenden

Sozialversicherungsträger sind gesonders geregelt.

Höhe des Schadenersatzes
- derzeitiger Betrag
- künftige personenbezogene Änderungen
- Mitverschuldensanteil/Quotenbildung dem Grunde nach
- Laufzeit
- Zinsfuss
- Zahlungsweise



Definieren Sie die grundsätzlichen Vorraussetzungen zur Deckung in den ProdHB 
Schäden durch
- Erbrachte Arbeiten oder sonstige Leistungen (+)
Zeitliche Abgrenzung: nach
• Abschluss der Arbeiten
• Ausführung der Leistung bzw.
• Inverkehrbringen des Erzeugnisses




Teilen Sie die Sektion 4 der ProdHB ein und kategorisieren sie die Weiterverarbeitungstatbestände

Die wesentlichen Termini sind:
Erzeugnis: Produkt, das VN geliefert hat
Produkt: Produkt, das ein Dritter zugeliefert hat und das zusammen mit dem Erzeugnis des VN verarbeitet wurde
Gesamtprodukt:das Ergebnis der Verarbeitung

4.2 Verbindung, Vermischung, Verarbeitung
4.3 Weiterbe- oder Weiterverarbeitung
4.4 Aus- und Einbau (+Einzelteileaustausch fak.)
4.5 Maschinenklausel (fak.)
4.6 Überprüfungskosten (fak.)

4.2: Weiterverarbeitung mit
Hinzufügen anderen Materials
4.2: Verbindung ist nicht trennbar (faktisch
oder wirtschaftlich)
4.3: Weiterbearbeitung ohne Hinzufügen anderen Materials
4.4: Verbindung ist trennbar

Es gelten das
Bausteinprinzip (aufgelistete
Verarbeitungssituationen als
Sachverhalte, an die die Deckung
anknüpft)
und
Enumerationsprinzip (nur die jeweils
genannten Schadenpositionen sind
gedeckt)


Differenzieren Sie Eigentumsverletzung und Sach/Vermögensschaden anhand des Transistorbeispiels

Zu Erinnerung
Transistoren Preis: 10 Cent
Weiterverarbeitung durch Dritten in Leiterplatte mit Preis: 5 €,
davon
2 € anderes Material,
2 € Herstellung (Verarbeitung),
90 Cent Gewinn
Diese werden wiederum in Steuergeräte verbaut.
Transistoren werden untrennbar mit Leiterplatten verbunden,
aus dem Steuergerät sind sie jedoch ausbaubar.
10.000 Einheiten sind betroffen.

Das Material ist eine Eigentumsverletzung § 823 Abs. 1 BGB
=> Haftung für 10.000 x 2 € = 20.000 €
=> Kosten für Transistoren nicht gedeckt! (4.2.2.2 ProdHB)

Der überschießender Wert Leiterplatte + Folgeschäden (Rückruf), sowie entgangener Gewinn:
und Herstellungskosten:
reiner Vermögensschaden
- nur vertragsrechtliche Haftung
- der Überschießender Wert der Leiterplatten war 3€, sprich 30,000
Tatsächlich sind diese alle gedeckt durch Klauseln der ProdHB (hätten es aber nicht alle sein müssen!).


Definieren Sie Nacherfüllung und die Anwendbarkeit in den ProdHB

Hat der Käufer die mangelhafte Sache gemäß ihrer Art und ihrem
Verwendungszweck in eine andere Sache eingebaut, oder an eine
andere Sache angebracht, ist der Verkäufer im Rahmen der
Nacherfüllung verpflichtet, dem Käufer die erforderlichen
Aufwendungen für das Entfernen der mangelhaften und den Einbau
oder das Anbringen der nachgebesserten oder gelieferten
mangelfreien Sache zu ersetzen.
(§ 439 BGB)

Jenachdem ob es sich um einen Einzelteil oder Erzeugnisaustausch handelt, sind weniger oder mehr
(genau in dieser Abfolge) Kosten gedeckt. 

Zu beachten ist trotzdem das allg. ein Erfüllungsschadenausschluss herrscht, d.h. Erfüllungssurrogate (substitutive Leistungen)
um den Auftrag zu Erfüllen sind nicht versichert.


Was sind die Grenzen beim Enumerationsprinzip bei Kosten der Überprüfung bei Mangelverdacht?

Beispielsweise muss nach Enumerationsprinzip folgendes getan werden bei Mangelverdacht bei Erzeugnissen
• notwendiges Vorsortieren zu
überprüfender Produkte
• Aussortieren von überprüften
Produkten
• infolge der Überprüfung
notwendiges Umpacken

Daher gilt die Abwägung in Ziff. 4.6.3: Begrenzung der gedeckten Kosten bei Überprüfung
Ggf. Überprüfung teurer als das Verwerfen oder Nachbearbeiten aller Teile, dann werden alle Produkte
so behandelt als wären sie mangelhaft.
